\documentclass[12pt,a4paper]{article}
\usepackage[utf8]{inputenc}
\usepackage[T1]{fontenc}
\usepackage[brazilian]{babel}
\usepackage{geometry}
\geometry{left=3cm,right=2cm,top=3cm,bottom=2cm}
\usepackage{setspace}
\onehalfspacing
\usepackage{array,tabularx,longtable,booktabs}
\usepackage{xcolor}
\usepackage{listings}
\usepackage[hidelinks]{hyperref}
\usepackage{enumitem}
\setlist{nosep}
\lstdefinelanguage{Pifari}{
  sensitive=true,
  morekeywords={if,else,elseif,for,while,break,continue,return,import,Link,to,agent,task,async,await,spawn,send,receive,close,try,catch,finally,throw,raise,match,case,const,dynamic,public,private,protected,static,override,extends,implements,as,from,in,where,reduce,map,filter},
  morekeywords=[2]{int,char,short,long,float,double,bool,void,signed,unsigned,byte,txt,Int,Float,Bool,Text,Null,None,List,Map,Set,Task,Agent}
}
\lstset{language=Pifari,basicstyle=\ttfamily\small,breaklines=true,columns=fullflexible,frame=single,showstringspaces=false}
\title{Definição da Linguagem de Programação Pifari\\Especificação Técnica — Versão 3}
\author{Documento de projeto}
\date{18 de outubro de 2026}
\begin{document}
\maketitle
\begin{abstract}
Este documento especifica a proposta da linguagem Pifari, com extensão \texttt{.pif}, orientada à composição de operações, programação concorrente e paralela, agentes nativos e interoperabilidade com C, Python e bibliotecas do sistema operacional. A versão estabelece uma base normativa de trabalho: sintaxe, léxico, tipos, semântica, fluxo, tarefas, memória, interoperabilidade, erros e conformidade. Onde uma escolha ainda exigir protótipo ou decisão de implementação, ela é identificada como tal, sem apresentar comportamento não implementado como fato.
\end{abstract}
\tableofcontents
\newpage
\section{Escopo, identidade e terminologia}
\subsection{Identidade}
Pifari é uma linguagem interpretada, de extensão \texttt{.pif}, concebida para equilibrar legibilidade, expressividade, segurança e acesso a recursos de baixo nível. A linguagem pretende oferecer interoperabilidade nativa com C e Python, abstrações de tarefas assíncronas e paralelas e agentes com estado próprio.
\subsection{Objetivos}
\begin{itemize}
\item interpretação direta do programa;
\item sistema de tipos híbrido: inferência estática sempre que possível e verificações dinâmicas quando necessárias;
\item tipos primitivos compatíveis conceitualmente com os tipos fundamentais de C, além de tipos de alto nível;
\item tarefas, agentes, mensagens e paralelismo coordenados pelo runtime;
\item memória gerenciada automaticamente, com recursos explícitos de baixo nível;
\item composição funcional e processamento de valores iteráveis;
\item interoperabilidade com bibliotecas C, módulos Python e bibliotecas do sistema operacional;
\item gramática e semântica especificadas antes da implementação de referência.
\end{itemize}
\subsection{Termos normativos}
As palavras \textbf{DEVE}, \textbf{NÃO DEVE}, \textbf{RECOMENDA-SE} e \textbf{PODE} indicam, respectivamente, requisito obrigatório, proibição, prática recomendada e recurso opcional. A especificação distingue: (a) decisões confirmadas pelo projeto; (b) regras recomendadas para a base normativa desta versão; (c) pontos cuja validação depende de protótipo. Uma implementação conforme deve documentar qualquer recurso ainda não suportado.
\subsection{Princípios}
O programa deve ser previsível, os efeitos colaterais devem ser identificáveis, as dependências de dados devem ser explícitas para o runtime e a concorrência não deve alterar silenciosamente a semântica sequencial. A otimização paralela só pode ocorrer quando a implementação preservar o comportamento observável definido pela linguagem.
\section{Modelo de execução}
\subsection{Interpretador}
A arquitetura de referência é um interpretador direto organizado em fases:
\begin{enumerate}
\item leitura UTF-8 e normalização de quebras de linha;
\item análise léxica e emissão de tokens com linha e coluna;
\item análise sintática e construção de AST;
\item resolução de nomes, escopos e imports;
\item análise de tipos e capacidades;
\item validação de efeitos, dependências e fronteiras de interoperabilidade;
\item execução da AST por uma máquina de avaliação;
\item escalonamento de tarefas, coleta de resultados e relatório de erros.
\end{enumerate}
O uso de AST é recomendado para diagnósticos e validação; uma implementação pode usar representação interna equivalente se preservar os requisitos observáveis.
\subsection{Runtime e encerramento}
O runtime gerencia o escalonamento sobre um conjunto de threads, a fila de tarefas, sincronização, cancelamento cooperativo, recursos externos e memória compartilhada. O programa não termina enquanto houver tarefas ou agentes pertencentes ao escopo do programa sem estado terminal. O encerramento global solicita o fechamento ordenado dos recursos e aguarda a finalização; tarefas bloqueadas em recursos externos devem ter política de timeout documentada.
\section{Tipos e valores}
\subsection{Tipos primitivos}
Pifari deve fornecer tipos de baixo nível compatíveis conceitualmente com C. O tamanho e o intervalo dos tipos de largura dependente da plataforma devem ser definidos pela implementação-alvo e expostos por constantes de plataforma; não se deve prometer largura universal quando C não a garante.
\begin{longtable}{>{\ttfamily}p{2.4cm}p{3.0cm}p{8.0cm}}
\toprule
Tipo Pifari & Correspondência conceitual & Significado\\
\midrule
void & void & ausência de valor de retorno; não é valor armazenável comum\\
bool & \texttt{\_Bool} / bool & valor lógico \texttt{true} ou \texttt{false}\\
char & char & unidade de caractere/byte conforme ABI; interoperabilidade exige convenção explícita\\
byte & unsigned char & unidade de 8 bits recomendada\\
short & short int & inteiro curto, com sinal por padrão\\
int & int & inteiro de tamanho definido pela ABI-alvo\\
long & long int & inteiro longo conforme ABI\\
long long & long long int & inteiro longo longo\\
unsigned short/int/long/long long & equivalentes C sem sinal & inteiros não negativos com largura da ABI\\
float & float & ponto flutuante de precisão simples\\
double & double & ponto flutuante de precisão dupla\\
long double & long double & precisão estendida conforme plataforma\\
\bottomrule
\end{longtable}
A grafia de tipos compostos por qualificadores deve ser aceita como sequência de tokens, por exemplo \texttt{unsigned long}. \texttt{char} não deve ser usado como sinônimo implícito de texto Unicode. Para texto, utiliza-se \texttt{txt} ou \texttt{Text}.
\subsection{Tipos de alto nível}
\begin{itemize}
\item \texttt{Text}: sequência Unicode;
\item \texttt{List[T]}: sequência ordenada de elementos;
\item \texttt{Map[K,V]}: associação entre chaves e valores;
\item \texttt{Set[T]}: coleção de valores únicos;
\item \texttt{Tuple[T1,\ldots,Tn]}: sequência de tamanho fixo;
\item \texttt{Range}: sequência numérica iterável;
\item \texttt{Option[T]}: valor presente ou ausente, se adotado pela biblioteca padrão;
\item \texttt{Result[T,E]}: resultado de sucesso ou erro, se adotado pela biblioteca padrão;
\item \texttt{Task[T]}: resultado futuro de uma operação;
\item \texttt{Agent}: entidade com estado próprio, mensagens e ciclo de vida;
\item \texttt{Null}: valor nulo da Pifari.
\end{itemize}
\subsection{Booleanos e ausência de valor}
\texttt{true} e \texttt{false} são literais booleanos. \texttt{Null} e \texttt{null} são grafias equivalentes para o valor nulo da Pifari; a implementação deve normalizar ambas para um único valor. \texttt{None} é reservado à interoperabilidade e representação de ausência oriunda do Python, salvo se uma biblioteca estabelecer conversão explícita. Não se deve presumir que \texttt{Null}, Python \texttt{None} e ponteiros C nulos sejam representações idênticas em memória.
\subsection{Literais textuais}
Aspas simples e duplas são aceitas: \texttt{"texto"} e \texttt{'texto'}. Para manter compatibilidade lexical com C e Python, escapes como \texttt{\textbackslash n}, \texttt{\textbackslash t}, \texttt{\textbackslash r}, \texttt{\textbackslash\textbackslash} e escape da aspa delimitadora devem ser definidos. Interpolação de texto, por exemplo \texttt{"Olá, \{nome\}"}, somente é válida quando o parser reconhece explicitamente a forma interpolada.
A forma \texttt{txt(conteúdo)} permite texto sem aspas. Recomendação normativa: o conteúdo sem aspas é texto literal até o parêntese de fechamento correspondente, não uma expressão. Parênteses aninhados, escapes e texto vazio devem ter regras próprias. Quando o texto contiver delimitadores ambíguos, deve-se usar aspas. A forma sem aspas não deve interpretar identificadores como variáveis.
\subsection{Coleções e inferência}
Coleções homogêneas são o padrão recomendado. Uma coleção heterogênea exige um tipo união explícito ou um tipo comum seguro; a implementação não deve apagar tipos silenciosamente. O tipo de uma variável inferida é obtido da inicialização. Inferência não equivale a permitir qualquer mudança posterior de tipo.
\section{Variáveis, atribuição e escopos}
\subsection{Declaração dinâmica e tipada}
\texttt{=} cria ou atualiza uma variável dinâmica, conforme as regras de ambiente e escopo. Exemplo: \texttt{idade = 10}. A forma \texttt{nome =: tipo(valor)} declara uma variável com tipo explícito e fixa seu tipo. Exemplo: \texttt{idade =: int(20)}. A atribuição posterior a uma variável tipada só é válida se o valor for compatível com o tipo declarado; não altera o tipo da variável.
\begin{lstlisting}
idade = 10
idade = "dez"             @ dinamica: tipo pode variar @
pontos =: int(20)
pontos = 21                @ valido: inteiro @
@ pontos = "vinte" invalido: tipo incompativel @
nome =: txt(Maria)
nome_c =: txt('Maria')
nome_python =: txt("Maria")
\end{lstlisting}
A declaração tipada exige construtor compatível com o tipo declarado. O cast não deve ser implícito: quando uma conversão for necessária, deve ser chamada uma função de conversão definida pela linguagem ou pelo programador.
\subsection{Mutabilidade, constantes e inicialização}
Uma variável comum é mutável. \texttt{const} declara um vínculo que não pode receber nova atribuição após inicializado. A mutabilidade do vínculo é distinta da mutabilidade interna do objeto. Toda variável deve ser inicializada antes de leitura, salvo se o tipo possuir valor padrão explicitamente especificado. Ler variável não declarada, não inicializada ou fora de escopo é erro.
\subsection{Escopo e visibilidade}
O escopo léxico é determinado por módulos, classes, funções e blocos. Reutilizar nome no mesmo escopo é proibido, salvo declaração de sobrecarga autorizada. Sombreamento em escopo interno deve ser documentado e diagnosticado. Atributos e métodos de classe são privados por padrão; acesso externo exige visibilidade explicitamente pública.
\section{Classes, atributos e métodos}
\subsection{Sintaxe declarativa}
\texttt{:=} inicia uma declaração de classe. \texttt{=:} introduz a seção de membros e declarações de atributos ou métodos, conforme a produção gramatical. Como \texttt{=:} também declara variáveis tipadas, o parser deve distinguir os usos pelo contexto: após o cabeçalho de classe, inicia membros; em uma instrução de atribuição, declara variável tipada. A distinção deve ser gramatical, nunca heurística.
\begin{lstlisting}
:= Pessoa
    =: nome txt
    =: idade int

    =: apresentar() 
        print(nome)
\end{lstlisting}
Este bloco ilustra a intenção, mas a forma definitiva dos cabeçalhos, construtores, parâmetros, tipos de retorno, inicializadores, herança e chamadas de método deve ser validada na gramática executável. Até que esses pontos sejam formalizados, o exemplo é não normativo.
\subsection{Contratos de classes}
A especificação completa deve definir: criação de instâncias; construtor; campos de instância e de classe; métodos de instância e estáticos; acesso privado/público/protegido; herança; sobrescrita; interfaces ou protocolos; inicialização; destruição e liberação de recursos; igualdade e representação textual. Nenhuma herança múltipla ou metaclasse é presumida nesta versão.
\section{Operadores e expressões}
\subsection{Inventário de operadores}
\begin{longtable}{p{3.1cm}p{3.3cm}p{6.8cm}}
\toprule
Operadores & Categoria & Regra\\
\midrule
\texttt{+ - * / // \% **} & aritméticos & soma, subtração, multiplicação, divisão real, divisão inteira, resto e potenciação\\
\texttt{== != < <= > >=} & comparação & igualdade, diferença e ordenação quando definida para o tipo\\
\texttt{and or not} & lógicos & conjunção, disjunção e negação booleana com curto-circuito\\
\texttt{=} & atribuição dinâmica & cria ou atualiza vínculo dinâmico\\
\texttt{=:} & declaração tipada/membros & declara tipo explícito; em classe introduz membros\\
\texttt{:=} & declaração de classe & inicia declaração de classe\\
\texttt{->} & lambda & introduz expressão lambda; não é operador geral de retorno\\
\texttt{|>} & fluxo direto & encadeia etapas; também pode iterar valores iteráveis conforme a forma gramatical\\
\texttt{<|} & aplicação/entrada & aplica o valor à operação à direita, conforme regra formal\\
\texttt{>|} & distribuição paralela & distribui trabalho entre ramos ou trabalhadores\\
\texttt{|<} & junção & sincroniza e reúne resultados dos ramos participantes\\
\texttt{.} & membro & acesso a membro público ou permitido\\
\texttt{[]} & indexação & acesso por índice/chave quando suportado\\
\texttt{()} & chamada/agrupamento & chamada de função e agrupamento\\
\bottomrule
\end{longtable}
\subsection{Precedência e associatividade recomendadas}
Da maior para a menor precedência:
\begin{enumerate}
\item agrupamento, acesso a membro, indexação e chamada;
\item operador unário \texttt{not}, sinal unário;
\item \texttt{**}, associativo à direita;
\item \texttt{* / // \%}, associativos à esquerda;
\item \texttt{+ -}, associativos à esquerda;
\item comparações \texttt{== != < <= > >=}; comparações encadeadas devem ser definidas ou proibidas explicitamente;
\item \texttt{and};
\item \texttt{or};
\item composição de fluxo \texttt{|> <|};
\item distribuição/junção \texttt{>| |<}, válidos apenas em contexto de fluxo;
\item atribuição e declarações, com restrições sintáticas próprias.
\end{enumerate}
A tabela é a base normativa recomendada e deve ser codificada em testes de parser. Parênteses são o mecanismo explícito para eliminar dúvidas. A atribuição não pode ser usada como expressão booleana por coerção implícita.
\subsection{Semântica aritmética}
\texttt{/} produz resultado de ponto flutuante. \texttt{//} calcula quociente inteiro conforme uma regra documentada para negativos; recomenda-se truncamento em direção a zero para alinhamento com C. \texttt{\%} deve satisfazer a relação entre dividendo, divisor e quociente escolhida para \texttt{//}. Divisão por zero gera erro de execução. Overflow de inteiros deve ser definido por modo: recomendado erro detectável, sem comportamento indefinido. \texttt{**} com expoente e base incompatíveis exige conversão explícita.
\subsection{Operadores lógicos e comparação}
Condições exigem \texttt{Bool}; números, textos e coleções não são convertidos implicitamente para booleanos. \texttt{and} e \texttt{or} usam curto-circuito. Comparações entre tipos diferentes são proibidas, salvo regra específica segura. Comparação de coleções deve ser definida por tipo; objetos sem relação de ordem não aceitam \texttt{<} ou \texttt{>}.
\section{Fluxo de controle}
\subsection{Condição direta}
A linguagem permite condição explícita \texttt{if}, \texttt{elseif} e \texttt{else}. Também prevê a forma concisa em que uma expressão booleana inicia um bloco:
\begin{lstlisting}
test = true
test
    print("ola")
else
    |> 0
\end{lstlisting}
A forma acima expressa a intenção de que um valor Bool possa iniciar uma decisão sem escrever \texttt{if}. A gramática deve definir como reconhecer o bloco, como associar \texttt{else} e se o resultado do bloco é um valor. Não se pode interpretar um valor não booleano dessa forma.
\subsection{Condicionais convencionais e valores}
\begin{lstlisting}
if (idade >= 18) {
    print("Maior de idade")
} elseif (idade >= 13) {
    print("Adolescente")
} else {
    print("Crianca")
}
\end{lstlisting}
A palavra reservada recomendada é \texttt{elseif}, como token único. Aceitar \texttt{else if} como sinônimo é opcional e deve ser especificado. Se uma condicional produzir valor, todos os caminhos que completam normalmente devem produzir valores compatíveis; um caminho que termina em erro ou não retorna pode ser tratado separadamente.
\subsection{Iteração}
A forma iterável desejada é:
\begin{lstlisting}
numeros |> numero |> print()
\end{lstlisting}
Interpretação recomendada: para cada elemento de \texttt{numeros}, vincular o elemento a \texttt{numero} e executar \texttt{print(numero)}. O nome de ligação é uma variável local ao corpo e não deve vazar para o escopo externo. Coleções ordenadas preservam sua ordem. Para conjuntos sem ordem definida, a ordem de iteração não é garantida.
\subsection{Laços explícitos}
\texttt{for} percorre um iterável; \texttt{while} repete enquanto a condição Bool for verdadeira; \texttt{break} encerra o laço mais interno; \texttt{continue} avança à próxima iteração; ambos são inválidos fora de laço. Laços não devem ser paralelizados automaticamente quando houver efeitos colaterais ou dependências não comprovadamente seguras.
\subsection{Capacidades}
A biblioteca e o sistema de tipos podem definir protocolos como \texttt{Iterable}, \texttt{Boolean}, \texttt{Comparable}, \texttt{Composable} e \texttt{Async}. A capacidade deve ser verificada estaticamente quando possível. A existência e os nomes finais desses protocolos precisam ser congelados antes da versão estável.
\section{Funções e módulos}
\subsection{Funções e lambdas}
\texttt{->} é reservado para expressões lambda. A definição de funções nomeadas deve ter uma forma única, com parâmetros, tipos opcionais, tipo de retorno opcional e corpo. Recomenda-se retorno explícito por \texttt{return}; retorno pela última expressão só deve ser habilitado se for uniforme em todos os contextos.
\begin{lstlisting}
dobrar = (x) -> x * 2
resultado = dobrar(5)
\end{lstlisting}
Este exemplo ilustra a lambda e não fixa a sintaxe definitiva de funções nomeadas. Parâmetros, argumentos nomeados, valores padrão, sobrecarga, recursão, funções genéricas e closures devem ter regras próprias.
\subsection{Módulos e importação}
\texttt{import} importa módulos Pifari ou módulos Python por um mecanismo de importação definido. A resolução de módulos deve ter ordem documentada, impedir ciclos perigosos ou detectá-los e informar caminhos de origem em erros. Importar um módulo não deve executar repetidamente sua inicialização, salvo especificação explícita.
\section{Concorrência, paralelismo e tarefas}
\subsection{Distinções}
Concorrência significa que tarefas têm progresso sobreposto; paralelismo significa execução simultânea em mais de um recurso. O runtime pode escalonar tarefas concorrentes sem garantir paralelismo físico em todo ambiente. A implementação não deve prometer desempenho linear nem execução paralela de operações bloqueadas pelo GIL do Python ou por limitações de biblioteca.
\subsection{Ciclo de vida da tarefa}
Estados conceituais: \texttt{Created}, \texttt{Ready}, \texttt{Running}, \texttt{Waiting}, \texttt{Succeeded}, \texttt{Failed} e \texttt{Cancelled}. Estados terminais são \texttt{Succeeded}, \texttt{Failed} e \texttt{Cancelled}. Uma tarefa criada não precisa começar imediatamente; deve ser iniciada pelo runtime ou por construção explícita definida na linguagem. A especificação deve escolher uma única regra consistente para o início. O resultado só pode ser lido após sucesso ou por operação assíncrona de espera.
\subsection{await e tarefas}
\texttt{await} espera a resolução de uma tarefa ou operação assíncrona e é válido apenas em contexto assíncrono ou em ponto de entrada que o runtime autorize. O uso de \texttt{await} em valor que não representa operação aguardável é erro de tipo. O resultado de uma tarefa falha propaga a falha ao ponto de espera, mas a falha permanece associada à tarefa original.
\subsection{Operadores de fluxo paralelo}
\begin{itemize}
\item \texttt{|>} encadeia etapas e, na forma iterável, aplica a operação a cada elemento;
\item \texttt{>|} inicia distribuição paralela explícita. Para coleções, a política padrão recomendada é dividir os elementos entre trabalhadores, preservando índice de origem;
\item \texttt{|<} é a barreira de junção: espera todas as tarefas do grupo e agrega os resultados na ordem original. Quando o operando à direita é um identificador ainda não ligado no escopo, ele recebe a lista agregada (forma de ligação); caso contrário, a expressão \texttt{|<} produz a lista;
\item \texttt{<|} avalia o fluxo à esquerda até a conclusão e só então avalia a expressão à direita; não inicia paralelismo por si só.
\end{itemize}
Nenhum operador de fluxo propaga escrita para o escopo do chamador: a variável de iteração e as variáveis locais vivem em um frame próprio do fluxo, e a base visível é somente leitura. Compartilhamento entre trabalhadores só ocorre por mecanismo explícito de célula protegida. O modo de distribuição (threads ou processos) e o dispositivo de execução não alteram o comportamento observável.
\subsection{Falhas e cancelamento}
A regra de projeto é que uma falha fica associada à tarefa que falhou. Tarefas independentes continuam; tarefas dependentes não podem consumir resultado inexistente e passam a falhar, aguardar recuperação ou ser canceladas conforme política explícita. O runtime não deve cancelar todo o processo por uma falha isolada, salvo se a falha for não recuperável ou o programa solicitar essa política.
O cancelamento é cooperativo. Tarefas verificam cancelamento em pontos seguros, como espera, início de etapa e chamadas de runtime. Operações nativas que não possam ser interrompidas devem declarar essa limitação. Recursos adquiridos devem ser liberados por mecanismo equivalente a \texttt{finally}/RAII.
\subsection{Agentes}
Agentes são construções nativas com estado próprio, capacidade de receber mensagens e executar tarefas assíncronas. Permanecem ativos até receber ordem de encerramento, falha fatal ou encerramento do programa. Mensagens e memória compartilhada são mecanismos igualmente importantes, mas não intercambiáveis: mensagens transferem dados segundo contrato; memória compartilhada exige proteção e sincronização.
A sintaxe final de declaração de agente, protocolo de mensagens, mailbox, backpressure, garantias de entrega, ordenação, supervisão e encerramento deve ser formalizada antes da implementação. Não se presume entrega exatamente uma vez para mensagens externas.
\subsection{Memória compartilhada}
O runtime protege acessos compartilhados por padrão. A implementação deve definir modelo de propriedade, referências, regiões compartilhadas, locks ou primitivas equivalentes, operações atômicas, visibilidade de escrita e prevenção de deadlocks. A segurança automática não pode significar que toda corrida lógica seja detectável. O programador deve poder solicitar mecanismos explícitos de baixo nível, com verificações e avisos adequados.
\section{Interoperabilidade}
\subsection{\texttt{import}}
\texttt{import} é destinado a módulos Pifari e módulos Python. Deve permitir alias e importação seletiva somente se essas formas forem especificadas na gramática. O adaptador Python deve definir versão suportada, gestão do interpretador Python, conversão de exceções e objetos, ciclo de vida de referências e limitações de concorrência.
\subsection{\texttt{\#include}}
\texttt{\#include} é uma diretiva de interoperabilidade e inclusão de declarações/cabeçalhos C e, mediante adaptadores específicos, interfaces C++ e C\#. A diretiva não deve fingir que cabeçalhos C, C++ e assemblies .NET têm a mesma semântica. A implementação deve encaminhar cada alvo ao adaptador correto, documentar a ferramenta de compilação/metadata e rejeitar combinações sem suporte.
\subsection{\texttt{Link to}}
\texttt{Link to} solicita ligação com bibliotecas ou símbolos disponibilizados pelo sistema operacional ou pelo próprio ecossistema Pifari. Pode referenciar bibliotecas dinâmicas, estáticas ou artefatos de ligação gerados pelo projeto. A resolução ocorre antes da execução do programa, na etapa de preparação/ligação; símbolos ausentes devem gerar erro de ligação claro. A gramática de caminhos, nomes de bibliotecas, símbolos e opções de ligação ainda deve ser definida.
\subsection{Compatibilidade de tipos e ABI}
A fronteira de interoperabilidade exige conversões explícitas ou declaradas em bindings. Inteiros devem respeitar largura e sinal; strings precisam de encoding, tamanho e propriedade de memória definidos; callbacks precisam de convenção de chamada e ciclo de vida documentados; ponteiros não devem ser tratados como objetos gerenciados sem wrapper. Conversões automáticas só podem ocorrer quando forem seguras, reversíveis ou expressamente configuradas.
\section{Memória e recursos}
A memória gerenciada é o padrão. O runtime pode utilizar contagem de referências, rastreamento, regiões ou estratégia híbrida, mas deve especificar ciclos, finalização e interação com objetos externos. Liberação de recursos externos não deve depender exclusivamente de coleta não determinística. Arquivos, sockets, locks, handles nativos e referências Python devem ter escopo e fechamento explícitos ou mecanismos determinísticos equivalentes.
\section{Regras léxicas}
\subsection{Codificação e identificadores}
O fonte usa UTF-8. Identificadores começam por letra ASCII ou sublinhado e continuam com letras ASCII, dígitos ou sublinhados; a linguagem diferencia maiúsculas e minúsculas. Recomenda-se não usar acentos em identificadores na primeira versão. Palavras reservadas não podem ser usadas como identificadores sem mecanismo de escape, se houver.
\subsection{Números}
Inteiros decimais usam dígitos, com sublinhado interno opcional se definido pela implementação. Reais usam ponto decimal, por exemplo \texttt{3.14}; notação científica deve seguir a forma \texttt{1.2e-3}. O sinal unário é operador, não parte do token numérico. Overflow em literal deve ser diagnosticado.
\subsection{Comentários}
Todo comentário começa com \texttt{@} e termina no próximo \texttt{@}. A regra vale para comentário curto ou multilinha. Não há aninhamento na versão inicial. Um \texttt{@} sem fechamento produz erro léxico apontando a posição de abertura. Como o delimitador pode aparecer em texto, dentro de literais textuais o lexer respeita primeiro o estado de string.
\begin{lstlisting}
@ comentario curto @
@ comentario longo
que pode atravessar varias linhas @
mensagem =: txt("O caractere @ dentro de string nao abre comentario")
\end{lstlisting}
\subsection{Espaços, tabulações e chaves}
Espaços e tabulações são permitidos. Quando a indentação for usada para legibilidade, ela não define o bloco por si só; chaves \texttt{\{} e \texttt{\}} delimitam blocos na sintaxe convencional. Para construções concisas de condição e iteração sem chaves, a gramática deve definir precisamente a extensão do bloco e a associação de \texttt{else}. Uma implementação não pode misturar regras de indentação obrigatória e chaves como se ambas fossem delimitadores independentes sem especificar sua interação. Recomenda-se que a primeira gramática normativa use chaves para blocos convencionais e trate formas sem chaves como produções separadas.
\subsection{Inventário de palavras reservadas}
A lista normativa inicial recomendada é:
\begin{longtable}{p{3.0cm}p{10.2cm}}
\toprule
Grupo & Palavras\\
\midrule
Lógicos e literais & \texttt{true}, \texttt{false}, \texttt{Null}, \texttt{null}\\
Controle & \texttt{if}, \texttt{else}, \texttt{elseif}, \texttt{for}, \texttt{while}, \texttt{break}, \texttt{continue}, \texttt{return}, \texttt{match}, \texttt{case}\\
Tarefas & \texttt{task}, \texttt{async}, \texttt{await}, \texttt{spawn}, \texttt{cancel}\\
Agentes e mensagens & \texttt{agent}, \texttt{send}, \texttt{receive}, \texttt{close}\\
Exceções & \texttt{try}, \texttt{catch}, \texttt{finally}, \texttt{throw}, \texttt{raise}\\
Declarações & \texttt{const}, \texttt{dynamic}, \texttt{public}, \texttt{private}, \texttt{protected}, \texttt{static}, \texttt{override}, \texttt{extends}, \texttt{implements}\\
Módulos e interoperabilidade & \texttt{import}, \texttt{from}, \texttt{as}, \texttt{Link}, \texttt{to}\\
Operações de coleção & \texttt{in}, \texttt{where}, \texttt{map}, \texttt{filter}, \texttt{reduce}\\
\bottomrule
\end{longtable}
Esta lista é uma reserva proposta; uma palavra só deve ser incluída no lexer normativo quando sua produção sintática estiver especificada. \texttt{int}, \texttt{char}, \texttt{float}, \texttt{double}, \texttt{bool}, \texttt{void}, \texttt{txt} e os demais nomes de tipos podem ser tratados como identificadores em contexto de tipo, reduzindo reservas desnecessárias.
\subsection{Inventário de pontuação e operadores}
\begin{longtable}{p{3.4cm}p{9.8cm}}
\toprule
Tokens & Uso\\
\midrule
\texttt{( ) [ ] \{ \} , ; . :} & agrupamento, coleção, bloco, separação, membro e anotação\\
\texttt{= =: := ->} & atribuição, declaração tipada/membro, declaração de classe e lambda\\
\texttt{+ - * / // \% **} & aritmética\\
\texttt{== != < <= > >=} & comparação\\
\texttt{and or not} & lógica booleana\\
\texttt{|> <| >| |<} & fluxo, aplicação, distribuição e junção\\
\texttt{\#include} & diretiva de interoperabilidade processada na etapa de preparação\\
\bottomrule
\end{longtable}
\section{Gramática livre de contexto (GLC)}
\subsection{Convenções}
A gramática abaixo descreve o núcleo de referência. \texttt{IDENT}, \texttt{INT}, \texttt{FLOAT}, \texttt{STRING}, \texttt{NEWLINE} e \texttt{EOF} são tokens lexicais. Comentários são removidos pelo lexer, mas as posições de origem são preservadas. A GLC não substitui as verificações semânticas de tipos, nomes, efeitos e dependências.
\begin{lstlisting}
programa       ::= lista_itens EOF
lista_itens    ::= item lista_itens | vazio
item           ::= declaracao | comando | expressao NEWLINE
declaracao     ::= declaracao_variavel
                 | declaracao_classe
                 | declaracao_funcao
                 | importacao
declaracao_variavel ::= IDENT "=" expressao
                      | IDENT "=:" tipo "(" expressao ")"
declaracao_classe ::= ":=" IDENT bloco_membros
bloco_membros  ::= NEWLINE INDENT lista_membros DEDENT
lista_membros  ::= membro lista_membros | vazio
membro         ::= "=:" IDENT tipo
                 | "=:" IDENT "(" parametros ")" bloco
declaracao_funcao ::= IDENT ":" "(" parametros ")" tipo_retorno bloco
parametros     ::= IDENT lista_parametros | vazio
lista_parametros ::= "," IDENT lista_parametros | vazio
tipo_retorno   ::= "->" tipo | vazio
tipo           ::= IDENT | tipo_generico
tipo_generico  ::= IDENT "[" tipo lista_tipos "]"
lista_tipos    ::= "," tipo lista_tipos | vazio
importacao     ::= "import" caminho_modulo
                 | "import" caminho_modulo "as" IDENT
bloco          ::= "{" lista_itens "}"
                 | NEWLINE INDENT lista_itens DEDENT
comando        ::= condicional | repeticao | retorno
                 | expressao NEWLINE
condicional    ::= "if" "(" expressao ")" bloco
                   lista_elseif opcional_else
                 | expressao_bool bloco lista_elseif opcional_else
lista_elseif   ::= "elseif" "(" expressao ")" bloco lista_elseif
                 | vazio
opcional_else  ::= "else" bloco | vazio
repeticao      ::= "while" "(" expressao ")" bloco
                 | "for" IDENT "in" expressao bloco
retorno        ::= "return" expressao_opcional NEWLINE
expressao_opcional ::= expressao | vazio
expressao      ::= atribuicao | fluxo
atribuicao     ::= IDENT "=" expressao
fluxo          ::= logica lista_fluxo
lista_fluxo    ::= operador_fluxo logica lista_fluxo | vazio
operador_fluxo ::= "|>" | "<|" | ">|" | "|<"
logica         ::= comparacao
                 | logica "and" comparacao
                 | logica "or" comparacao
                 | "not" logica
comparacao     ::= soma operador_comparacao soma
                 | soma
operador_comparacao ::= "==" | "!=" | "<" | "<=" | ">" | ">="
soma           ::= produto
                 | soma "+" produto
                 | soma "-" produto
produto        ::= potencia
                 | produto "*" potencia
                 | produto "/" potencia
                 | produto "//" potencia
                 | produto "%" potencia
potencia       ::= unario "**" potencia | unario
unario         ::= "-" primario | "+" primario | primario
primario       ::= literal | IDENT | chamada | "(" expressao ")"
                 | lista | lambda
chamada        ::= IDENT "(" argumentos ")"
argumentos     ::= expressao lista_argumentos | vazio
lista_argumentos ::= "," expressao lista_argumentos | vazio
lista          ::= "[" argumentos "]"
lambda         ::= IDENT "->" expressao
literal        ::= INT | FLOAT | STRING | "true" | "false" | "Null"
caminho_modulo ::= IDENT
\end{lstlisting}
\subsection{Limites da gramática de referência}
A GLC acima é um esqueleto normativo inicial, não uma gramática de produção completa. Antes do parser estável, devem ser corrigidos e congelados: lista de itens e terminadores; declarações tipadas versus atribuição; tipos C compostos; texto sem aspas em \texttt{txt}; expressões condicionais; blocos implícitos; classes; sintaxe dos ramos paralelos; precedência dos quatro operadores de fluxo; diretivas de pré-processamento; e recuperação de erro. A gramática final deve ser não ambígua ou declarar explicitamente as regras de desambiguação.
\section{Semântica estática e dinâmica}
\subsection{Análise estática}
Antes da execução, o interpretador deve resolver nomes e verificar, sempre que possível: declarações duplicadas, variáveis não declaradas, inicialização, compatibilidade de atribuições tipadas, argumentos de chamadas, retorno de função, uso de \texttt{break}/\texttt{continue}, uso de \texttt{await}, capacidades iteráveis e compatibilidade entre ramos de fluxo. Erros que dependem de valores dinâmicos podem ser detectados em execução.
\subsection{Conversões}
Todas as conversões entre tipos são explícitas, salvo promoção aritmética segura que seja formalmente aprovada. Conversões nativas e funções definidas pelo programador são permitidas. Cada conversão deve definir: tipos de origem e destino; comportamento em overflow; perda de precisão; falha de conversão; e efeito sobre propriedade e memória. Não existe cast universal que torne qualquer valor válido sem verificação.
\subsection{Efeitos e determinismo}
Funções devem ter efeitos observáveis classificáveis: pura, leitura, escrita, I/O, tarefa, agente ou chamada externa. A análise de efeitos pode impedir paralelização automática de operações incompatíveis. Quando o runtime não consegue provar segurança, deve executar sequencialmente ou exigir anotação explícita, em vez de assumir independência.
\subsection{Semântica operacional resumida}
A avaliação de uma expressão produz um valor, uma falha ou uma tarefa pendente. Uma instrução pode alterar o ambiente local, criar tarefa ou realizar efeito externo. A sequência de instruções preserva ordem observável. Para um fluxo paralelo, o runtime registra o grupo de tarefas, dependências e índices de origem; a junção aguarda os participantes e agrega resultados conforme política declarada. Falhas permanecem vinculadas às tarefas e são expostas ao ponto de junção ou espera.
\section{Erros e diagnósticos}
A implementação deve distinguir:
\begin{itemize}
\item erro léxico: caractere inválido, string ou comentário sem fechamento, literal numérico malformado;
\item erro sintático: token inesperado, delimitador ausente, bloco inválido ou sequência de operadores sem operandos;
\item erro semântico: nome inexistente, declaração duplicada, tipo incompativel, retorno inválido ou uso indevido de capacidade;
\item erro de execução: divisão por zero, cast inválido, índice fora dos limites, falha externa ou exceção não tratada;
\item erro de ligação: biblioteca, símbolo ou ABI indisponível;
\item erro de concorrência: deadlock detectado, tarefa cancelada, dependência falha ou uso proibido de recurso compartilhado.
\end{itemize}
Toda mensagem deve conter categoria, mensagem legível, arquivo, linha e coluna quando disponíveis, trecho de código e sugestão de correção quando segura. Uma falha de tarefa deve preservar a pilha de chamadas e o identificador da tarefa. O runtime deve evitar continuar em estado inconsistente.
\section{Testes de conformidade}
Cada regra normativa deve ter testes positivos e negativos. A suíte mínima deve conter:
\begin{enumerate}
\item lexer: comentários curtos/multilinha, aspas simples/duplas, escapes, UTF-8, tokens compostos e operadores;
\item parser: precedência, associatividade, blocos, classes, declarações e erros de delimitadores;
\item tipos: inferência, \texttt{=:}, variáveis dinâmicas, casts e compatibilidade de tipos C;
\item controle: \texttt{if}/\texttt{elseif}/\texttt{else}, condição booleana direta, iteração e escopo;
\item fluxo: cada operador isolado, composições, ordem, erros e expressões entre parênteses;
\item concorrência: dependências, resultados ordenados, falhas isoladas, cancelamento, encerramento e ausência de deadlock nos casos de teste;
\item interoperabilidade: cabeçalho ausente, símbolo ausente, ABI incompatível, conversão de strings, memória e exceções Python;
\item regressão: todo bug corrigido gera caso reproduzível.
\end{enumerate}
A especificação só pode ser considerada pronta para implementação estável quando cada produção gramatical e cada regra semântica tiver teste associado ou justificativa explícita.
\section{Exemplos integrados}
\subsection{Tipos, textos e comentários}
\begin{lstlisting}
@ Exemplo de declaracao tipada @
idade =: int(20)
nome =: txt(Maria)
saudacao =: txt("Ola, " + nome)
ativo = true
print(saudacao)
\end{lstlisting}
A expressão de concatenação acima depende de uma operação de texto explicitamente definida; a versão normativa deve escolher entre concatenação por \texttt{+} para dois Text ou função própria. A concatenação de texto e número não é implícita.
\subsection{Iteração por fluxo}
\begin{lstlisting}
numeros = [1, 2, 3, 4, 5]
numeros |> numero |> print(numero)
\end{lstlisting}
Este exemplo representa a forma concisa desejada. A gramática final deve definir se a etapa final recebe o elemento como argumento implícito ou exige que a variável seja passada explicitamente; os dois modelos não podem coexistir sem regra clara.
\subsection{Paralelismo conceitual}
\begin{lstlisting}
dados = [1, 2, 3, 4, 5, 6]
resultados = dados >| processar |<
print(resultados)
\end{lstlisting}
Este exemplo é ilustrativo. A sintaxe final deve especificar a função de processamento, a forma de declarar o grupo de trabalhadores e como o resultado da junção é atribuído. A semântica recomendada preserva a ordem original e aguarda todas as tarefas participantes.
\subsection{Interoperabilidade conceitual}
\begin{lstlisting}
#include <stdio.h>
Link to "biblioteca_sistema"
import modulo_python
\end{lstlisting}
O exemplo demonstra as três intenções distintas. A sintaxe concreta da diretiva \texttt{\#include}, a resolução do nome após \texttt{Link to} e a integração de módulos Python dependem da implementação e do alvo de plataforma.
\section{Plano de implementação após congelamento}
A implementação deve seguir uma sequência que minimize retrabalho:
\begin{enumerate}
\item congelar este documento após revisão das ambiguidades explicitamente listadas;
\item construir lexer independente, com testes de tokens e posições;
\item implementar parser da GLC e testes de árvore sintática;
\item implementar resolução de nomes, escopos e tipos;
\item implementar avaliador sequencial de expressões e controle de fluxo;
\item implementar classes, módulos e conversões;
\item implementar runtime de tarefas, estados, cancelamento e junção;
\item adicionar agentes, mensagens e memória compartilhada;
\item implementar adaptadores de interoperabilidade C, Python e bibliotecas de sistema;
\item executar testes de conformidade, benchmarks e testes de estresse;
\item publicar uma matriz de recursos suportados por plataforma.
\end{enumerate}
Não se deve começar pela implementação completa dos operadores paralelos antes de congelar a semântica de valores, efeitos, falhas e dependências. Um protótipo pode explorar opções, mas não deve ser tratado como referência normativa.
\section{Questões que exigem validação final}
Antes de declarar a gramática estável, a equipe deve decidir formalmente:
\begin{enumerate}
\item se blocos convencionais usarão chaves, indentação ou ambos sob regras distintas;
\item sintaxe completa de classes, construtores, membros, herança e tipos de retorno;
\item semântica exata de \texttt{<|}, \texttt{>|} e \texttt{|<}, incluindo a forma concreta de ramificação;
\item se \texttt{txt(conteúdo sem aspas)} aceita espaços, parênteses aninhados e escapes e como encerra o literal;
\item se a iteração \texttt{colecao |> elemento |> operação} é uma produção dedicada ou açúcar sintático;
\item largura e overflow de tipos numéricos por plataforma;
\item política de propriedade, cópia e compartilhamento de objetos entre tarefas;
\item versão Python, ABI C, ferramentas de ligação e estratégia de suporte a C++/.NET;
\item formato e escopo de \texttt{\#include} e \texttt{Link to};
\item modelo definitivo de exceções, falhas de tarefa e agentes.
\end{enumerate}
\section{Conclusão}
A Pifari tem como eixo a combinação de expressividade, tipos, interoperabilidade e execução concorrente administrada pelo runtime. A qualidade da linguagem dependerá menos da quantidade de símbolos e mais da consistência entre léxico, gramática, tipos, semântica e testes. Esta versão oferece uma especificação de trabalho ampliada e identifica explicitamente o que ainda precisa de validação. Ela não afirma que o interpretador, o runtime ou os adaptadores já existam.
\end{document}
